% Options for packages loaded elsewhere
\PassOptionsToPackage{unicode}{hyperref}
\PassOptionsToPackage{hyphens}{url}
\PassOptionsToPackage{dvipsnames,svgnames,x11names}{xcolor}
%
\documentclass[
]{scrartcl}

\usepackage{amsmath,amssymb}
\usepackage{iftex}
\ifPDFTeX
  \usepackage[T1]{fontenc}
  \usepackage[utf8]{inputenc}
  \usepackage{textcomp} % provide euro and other symbols
\else % if luatex or xetex
  \usepackage{unicode-math}
  \defaultfontfeatures{Scale=MatchLowercase}
  \defaultfontfeatures[\rmfamily]{Ligatures=TeX,Scale=1}
\fi
\usepackage{lmodern}
\ifPDFTeX\else  
    % xetex/luatex font selection
\fi
% Use upquote if available, for straight quotes in verbatim environments
\IfFileExists{upquote.sty}{\usepackage{upquote}}{}
\IfFileExists{microtype.sty}{% use microtype if available
  \usepackage[]{microtype}
  \UseMicrotypeSet[protrusion]{basicmath} % disable protrusion for tt fonts
}{}
\makeatletter
\@ifundefined{KOMAClassName}{% if non-KOMA class
  \IfFileExists{parskip.sty}{%
    \usepackage{parskip}
  }{% else
    \setlength{\parindent}{0pt}
    \setlength{\parskip}{6pt plus 2pt minus 1pt}}
}{% if KOMA class
  \KOMAoptions{parskip=half}}
\makeatother
\usepackage{xcolor}
\setlength{\emergencystretch}{3em} % prevent overfull lines
\setcounter{secnumdepth}{5}
% Make \paragraph and \subparagraph free-standing
\ifx\paragraph\undefined\else
  \let\oldparagraph\paragraph
  \renewcommand{\paragraph}[1]{\oldparagraph{#1}\mbox{}}
\fi
\ifx\subparagraph\undefined\else
  \let\oldsubparagraph\subparagraph
  \renewcommand{\subparagraph}[1]{\oldsubparagraph{#1}\mbox{}}
\fi

% Lower the position of page numbers
\usepackage{fancyhdr}
\fancyhf{}
\fancyfoot[C]{\thepage}
\setlength{\footskip}{100pt} % Increase this value to lower the page numbers



\providecommand{\tightlist}{%
  \setlength{\itemsep}{0pt}\setlength{\parskip}{0pt}}\usepackage{longtable,booktabs,array}
\usepackage{calc} % for calculating minipage widths
% Correct order of tables after \paragraph or \subparagraph
\usepackage{etoolbox}
\makeatletter
\patchcmd\longtable{\par}{\if@noskipsec\mbox{}\fi\par}{}{}
\makeatother
% Allow footnotes in longtable head/foot
\IfFileExists{footnotehyper.sty}{\usepackage{footnotehyper}}{\usepackage{footnote}}
\makesavenoteenv{longtable}
\usepackage{graphicx}
\makeatletter
\newsavebox\pandoc@box
\newcommand*\pandocbounded[1]{% scales image to fit in text height/width
  \sbox\pandoc@box{#1}%
  \Gscale@div\@tempa{\textheight}{\dimexpr\ht\pandoc@box+\dp\pandoc@box\relax}%
  \Gscale@div\@tempb{\linewidth}{\wd\pandoc@box}%
  \ifdim\@tempb\p@<\@tempa\p@\let\@tempa\@tempb\fi% select the smaller of both
  \ifdim\@tempa\p@<\p@\scalebox{\@tempa}{\usebox\pandoc@box}%
  \else\usebox{\pandoc@box}%
  \fi%
}
% Set default figure placement to htbp
\def\fps@figure{htbp}
\makeatother

% Change font size
\usepackage[font={footnotesize, sf}, labelfont=bf, margin=0.05\linewidth]{caption}
\renewcommand{\familydefault}{\sfdefault}
% Bold the 'Figure #' in the caption and separate it from the title/caption
% with a period
% Captions will be left justified
\usepackage[
	aboveskip=1pt,
	labelfont=bf,
	labelsep=period,
	justification=raggedright,
	singlelinecheck=off
]{caption}
\renewcommand{\figurename}{Fig}

% Add numbered lines
\usepackage{lineno}
% \linenumbers

% Package to include multiple title pages
% This will allow me to add a tile to the main text and to the SI
\usepackage{titling}

% Define command to begin the supplementary section
\newcommand{\beginsupplement}{
\setcounter{page}{1} % Restart page counter
\setcounter{section}{0} % Restart section counter
% \renewcommand{\thesection}{\Alph{section}}%
\renewcommand{\thesection}{}
\renewcommand{\thesubsection}{\Alph{subsection}}
\setcounter{table}{0} % Restart table counter
\renewcommand{\thetable}{\Alph{table}}
\setcounter{figure}{0} % Restart figure counter
\renewcommand{\thefigure}{\Alph{figure}}%
\setcounter{equation}{0} % Restart equation counter
\renewcommand{\theequation}{S\arabic{equation}}%
}

% Add author affiliations
\usepackage{authblk}

% Allow to define personalized colors
\usepackage{xcolor}
\makeatletter
\@ifpackageloaded{caption}{}{\usepackage{caption}}
\AtBeginDocument{%
\ifdefined\contentsname
  \renewcommand*\contentsname{Table of contents}
\else
  \newcommand\contentsname{Table of contents}
\fi
\ifdefined\listfigurename
  \renewcommand*\listfigurename{List of Figures}
\else
  \newcommand\listfigurename{List of Figures}
\fi
\ifdefined\listtablename
  \renewcommand*\listtablename{List of Tables}
\else
  \newcommand\listtablename{List of Tables}
\fi
\ifdefined\figurename
  \renewcommand*\figurename{Figure}
\else
  \newcommand\figurename{Figure}
\fi
\ifdefined\tablename
  \renewcommand*\tablename{Table}
\else
  \newcommand\tablename{Table}
\fi
}
\@ifpackageloaded{float}{}{\usepackage{float}}
\floatstyle{ruled}
\@ifundefined{c@chapter}{\newfloat{codelisting}{h}{lop}}{\newfloat{codelisting}{h}{lop}[chapter]}
\floatname{codelisting}{Listing}
\newcommand*\listoflistings{\listof{codelisting}{List of Listings}}
\makeatother
\makeatletter
\makeatother
\makeatletter
\@ifpackageloaded{caption}{}{\usepackage{caption}}
\@ifpackageloaded{subcaption}{}{\usepackage{subcaption}}
\makeatother
\ifLuaTeX
  \usepackage{selnolig}  % disable illegal ligatures
\fi

%%%%%%%%%%%%%%%%%%%%%%%%%%%%%%%%%%%%%%%%%%%%%%%%%%%%%%%%%%%%%%%%%%%%%%%%%%%%%%%%
\usepackage[
    style=vancouver,
    sorting=none,				% Do not sort bibliography
	url=false, 					% Do not show url in reference
	doi=false, 					% Do not show doi in reference
	isbn=false, 				% Do not show isbn link in reference
	eprint=false, 			    % Do not show eprint link in reference
	maxbibnames=9, 			    % Include up to 9 names in citation
	firstinits=true,
]{biblatex}
%%%%%%%%%%%%%%%%%%%%%%%%%%%%%%%%%%%%%%%%%%%%%%%%%%%%%%%%%%%%%%%%%%%%%%%%%%%%%%%%
\addbibresource{references.bib}
\IfFileExists{bookmark.sty}{\usepackage{bookmark}}{\usepackage{hyperref}}
\IfFileExists{xurl.sty}{\usepackage{xurl}}{} % add URL line breaks if available
\urlstyle{same} % disable monospaced font for URLs
\hypersetup{
  pdftitle={Code Overview},
  pdfauthor={Manuel Razo-Mejia},
  pdfkeywords={variational autoencoders, evolutionary
dynamics, antibiotic resistance},
  colorlinks=true,
  linkcolor={blue},
  filecolor={Maroon},
  citecolor={Blue},
  urlcolor={Blue},
  pdfcreator={LaTeX via pandoc}}

%%%%%%%%%%%%%%%%%%%%%%%%%%%%%%%%%%%%%%%%%%%%%%%%%%%%%%%%%%%%%%%%%%%%%%%%%%%%%%%%
% Add title
\title{Code Overview}
%%%%%%%%%%%%%%%%%%%%%%%%%%%%%%%%%%%%%%%%%%%%%%%%%%%%%%%%%%%%%%%%%%%%%%%%%%%%%%%%

%%%%%%%%%%%%%%%%%%%%%%%%%%%%%%%%%%%%%%%%%%%%%%%%%%%%%%%%%%%%%%%%%%%%%%%%%%%%%%%%
% Add list of authors

% Loop through authors
\author[1]{Manuel Razo-Mejia}

% Loop through affiliations
\affil[1]{Department of Biology, Stanford University}

% Add corresponding author
%%%%%%%%%%%%%%%%%%%%%%%%%%%%%%%%%%%%%%%%%%%%%%%%%%%%%%%%%%%%%%%%%%%%%%%%%%%%%%%%

%%%%%%%%%%%%%%%%%%%%%%%%%%%%%%%%%%%%%%%%%%%%%%%%%%%%%%%%%%%%%%%%%%%%%%%%%%%%%%%%
% Set affiliations in small font
\renewcommand\Affilfont{\itshape\small}
%%%%%%%%%%%%%%%%%%%%%%%%%%%%%%%%%%%%%%%%%%%%%%%%%%%%%%%%%%%%%%%%%%%%%%%%%%%%%%%%

%%%%%%%%%%%%%%%%%%%%%%%%%%%%%%%%%%%%%%%%%%%%%%%%%%%%%%%%%%%%%%%%%%%%%%%%%%%%%%%%
% Remove date
\date{}
%%%%%%%%%%%%%%%%%%%%%%%%%%%%%%%%%%%%%%%%%%%%%%%%%%%%%%%%%%%%%%%%%%%%%%%%%%%%%%%%

% Add package for line numbering
\usepackage{lineno}

\begin{document}
\linenumbers
\maketitle
\section{Code Documentation}\label{code-documentation}

This section provides interactive code examples and demonstrations of
the computational methods used in our research on evolutionary
landscapes and phenotype-to-fitness maps.

\subsection{Available Notebooks}\label{available-notebooks}

We provide several computational notebooks that showcase different
aspects of our methodology and analysis:

\begin{enumerate}
\def\labelenumi{\arabic{enumi}.}
\tightlist
\item
  \textbf{\href{code/evolutionary_dynamics.qmd}{Evolutionary Dynamics
  Simulation}}
\end{enumerate}

\begin{itemize}
\tightlist
\item
  Simulates the dynamics of populations evolving in a 2-dimensional
  phenotypic space under various fitness landscapes and fixed mutational
  constraints. This notebook demonstrates how populations navigate
  phenotypic space influenced by both fitness gradients and
  genotype-phenotype density.
\end{itemize}

\begin{enumerate}
\def\labelenumi{\arabic{enumi}.}
\setcounter{enumi}{1}
\tightlist
\item
  \textbf{\href{code/rhvae_sim.qmd}{RHVAE on Simulated Data}}
\end{enumerate}

\begin{itemize}
\tightlist
\item
  Implements a Riemannian Hamiltonian Variational Autoencoder (RHVAE) to
  analyze fitness profiles generated from evolutionary simulations. This
  notebook shows how we can learn low-dimensional representations of
  high-dimensional fitness data while preserving the geometric structure
  of the underlying phenotypic space.
\end{itemize}

\begin{enumerate}
\def\labelenumi{\arabic{enumi}.}
\setcounter{enumi}{2}
\tightlist
\item
  \textbf{\href{code/iwasawa_bayesian_inference.qmd}{Bayesian Inference
  of IC₅₀ Values}}
\end{enumerate}

\begin{itemize}
\tightlist
\item
  Performs Bayesian inference on the IC₅₀ values of antibiotic
  resistance landscapes using experimental data. This notebook
  demonstrates model fitting of dose-response curves using different
  parameterizations and implements outlier detection techniques.
\end{itemize}

\begin{enumerate}
\def\labelenumi{\arabic{enumi}.}
\setcounter{enumi}{3}
\tightlist
\item
  \textbf{\href{code/metropolis_kimura_evolution.qmd}{Metropolis-Kimura
  Evolution}}
\end{enumerate}

\begin{itemize}
\tightlist
\item
  Explores a biologically motivated model of evolution combining
  elements of the Metropolis-Hastings algorithm with Motoo Kimura's
  population genetics theory. This model accounts separately for
  mutation accessibility and selection via fixation probability.
\end{itemize}

\begin{enumerate}
\def\labelenumi{\arabic{enumi}.}
\setcounter{enumi}{4}
\tightlist
\item
  \textbf{\href{code/metropolis_hastings_evolution.qmd}{Metropolis-Hastings
  Evolution}}
\end{enumerate}

\begin{itemize}
\tightlist
\item
  Implements the Metropolis-Hastings algorithm for simulating
  evolutionary dynamics on fitness landscapes with mutational
  constraints.
\end{itemize}

\subsection{Running the Code}\label{running-the-code}

All notebooks can be run interactively if you have Julia installed with
the required packages. Each notebook is self-contained with detailed
explanations, mathematical background, and visualization of results.

\subsubsection{Required Packages}\label{required-packages}

To run these notebooks, you'll need:

\begin{itemize}
\tightlist
\item
  \textbf{Data Processing}: \texttt{DataFrames}, \texttt{CSV},
  \texttt{DimensionalData}
\item
  \textbf{Scientific Computing}: \texttt{LinearAlgebra},
  \texttt{StatsBase}, \texttt{Distributions}, \texttt{LsqFit}
\item
  \textbf{Bayesian Inference}: \texttt{Turing}
\item
  \textbf{Deep Learning}: \texttt{Flux}, \texttt{AutoEncoderToolkit}
\item
  \textbf{Visualization}: \texttt{CairoMakie}, \texttt{ColorSchemes}
\end{itemize}

Click on any of the notebooks in the sidebar to explore the code in
detail.



\end{document}
